\documentclass[12pt,reqno]{amsart}

\usepackage[T1]{fontenc}
\usepackage{lmodern}
\usepackage{microtype}
\usepackage{mathtools}
\usepackage{amssymb}
\usepackage{booktabs}
\usepackage{enumitem}
\usepackage{xcolor}
\usepackage[colorlinks=true,
  linkcolor=blue!55!black,
  citecolor=green!40!black,
  urlcolor=blue!60!black]{hyperref}
\hypersetup{
  pdftitle={Degree-sixteen cohomology of moduli spaces of stable curves in genus at most seven is Tate},
  pdfauthor={GPT-5.6, Claude Opus 5, Qwen 3.8, Grok 4.6, and Daisuke Ken Sakai},
  pdfsubject={Rational Hodge structures on moduli spaces of stable curves},
  pdfkeywords={moduli of curves, mixed Hodge structures, Tate Hodge structures, boundary stratification, tautological classes}}

\newcommand{\Q}{\mathbf{Q}}
\newcommand{\C}{\mathbf{C}}
\newcommand{\M}{\mathcal{M}}
\newcommand{\Mbar}{\overline{\mathcal{M}}}
\newcommand{\BM}{\mathrm{BM}}
\newcommand{\RH}{\mathrm{RH}}
\newcommand{\Mod}{\operatorname{Mod}}
\newcommand{\PMod}{\operatorname{PMod}}
\newcommand{\Aut}{\operatorname{Aut}}
\newcommand{\vcd}{\operatorname{vcd}}
\newcommand{\im}{\operatorname{im}}

\newtheorem{theorem}{Theorem}[section]
\newtheorem{proposition}[theorem]{Proposition}
\newtheorem{lemma}[theorem]{Lemma}
\newtheorem{corollary}[theorem]{Corollary}
\theoremstyle{definition}
\newtheorem{definition}[theorem]{Definition}
\theoremstyle{remark}
\newtheorem{remark}[theorem]{Remark}

\title[Degree-sixteen cohomology in genus at most seven]
  {Degree-sixteen cohomology of moduli spaces of stable curves in genus at
   most seven is Tate}
\author[GPT-5.6 et al.]
  {GPT-5.6, Claude Opus 5, Qwen 3.8, Grok 4.6,\\
   and Daisuke Ken Sakai}
\date{2 September 2026}
\subjclass[2020]{Primary 14H10; Secondary 14C30}
\keywords{moduli of curves, mixed Hodge structures, Tate Hodge structures,
boundary stratification, tautological classes}

\begin{document}

\begin{abstract}
We prove that, for every stable pair $(g,n)$ with $g\leq7$, the rational
Hodge structure $H^{16}(\overline{\mathcal M}_{g,n};\Q)$ is a direct sum of
copies of $\Q(-8)$.  The proof extends the boundary induction of
Canning--Larson--Payne by one even degree.  Its new open inputs are obtained
from their primitive-quotient formula, virtual cohomological dimension, and
the Ionel--Looijenga vanishing theorem for tautological classes.  The open Borel--Moore
term vanishes whenever $3\leq g\leq7$ and $2g+n\geq11$.  In particular, the
lowest-weight contribution from $\mathcal M_{6,6}$ is zero, so boundary
homology surjects onto $H_{16}(\overline{\mathcal M}_{6,6};\Q)$.  We also
show that the unrestricted degree-sixteen problem reduces to two open
inputs, on $\mathcal M_{8}$ and $\mathcal M_{8,1}$.
\end{abstract}

\maketitle

\section{Introduction}

Let $\Mbar_{g,n}$ be the smooth proper Deligne--Mumford stack of stable
complex curves of genus $g$ with $n$ ordered marked points, and let
$\M_{g,n}\subset\Mbar_{g,n}$ be the open substack of smooth pointed curves.
A pair $(g,n)$ is called stable if $2g-2+n>0$.  All cohomology and homology
groups in this paper have rational coefficients and carry their natural
mixed Hodge structures.

Canning, Larson, and Payne proved that the rational cohomology of
$\Mbar_{g,n}$ is Tate in every even degree at most fourteen, uniformly in
$g$ and $n$ \cite[Theorem~1.5(3) and the discussion following it]{CLP}.
Degree sixteen is the first even degree not covered by this result.  Its
status in the general Hodge- and Tate-theoretic predictions is discussed by
Payne \cite[Section~4]{PayneSurvey}.  We prove the prediction through genus
seven.

\begin{theorem}\label{thm:main}
Let $(g,n)$ be stable and suppose that $g\leq7$.  If
\[
  b_{g,n}:=\dim_{\Q}H^{16}(\Mbar_{g,n};\Q),
\]
then there is an isomorphism of polarizable rational Hodge structures
\begin{equation}\label{eq:main}
  H^{16}(\Mbar_{g,n};\Q)
  \cong \Q(-8)^{\oplus b_{g,n}}.
\end{equation}
\end{theorem}

The assertion is about Hodge type.  It neither determines $b_{g,n}$ nor
asserts that the group in \eqref{eq:main} is generated by algebraic or
tautological cycle classes.  This distinction is essential: Tate Hodge type
alone does not imply the Hodge conjecture.

The case that motivated the argument has an additional conclusion.  Write
$\widetilde{\partial\M}_{g,n}$ for the normalization of the reduced boundary
$\Mbar_{g,n}\setminus\M_{g,n}$.

\begin{corollary}\label{cor:66-intro}
The proper pushforward
\begin{equation}\label{eq:66-boundary-surj-intro}
 H_{16}(\widetilde{\partial\M}_{6,6};\Q)
 \longrightarrow H_{16}(\Mbar_{6,6};\Q)
\end{equation}
is surjective.  In particular, if
$b_{16}=\dim_{\Q}H^{16}(\Mbar_{6,6};\Q)$, then
\[
 H^{16}(\Mbar_{6,6};\Q)\cong\Q(-8)^{\oplus b_{16}}.
\]
\end{corollary}

The surjectivity follows from a vanishing on the open moduli stack:
\begin{equation}\label{eq:66-open-intro}
 W_{26}H^{26}(\M_{6,6};\Q)=0,
 \qquad
 W_{-16}H^{\BM}_{16}(\M_{6,6};\Q)=0.
\end{equation}
More generally, the critical open groups that occur at the boundary of the
known Chow--K\"unneth ranges all vanish:
\begin{equation}\label{eq:critical-intro}
\begin{split}
 W_{20}H^{20}(\M_{3,12};\Q)&=0,\\
 W_{22}H^{22}(\M_{4,10};\Q)&=0,\\
 W_{24}H^{24}(\M_{5,8};\Q)&=0,\\
 W_{26}H^{26}(\M_{6,6};\Q)&=0,\\
 W_{28}H^{28}(\M_{7,4};\Q)&=0.
\end{split}
\end{equation}
The second row lies in the published Chow--K\"unneth range.  The other four
rows, including the genus-five case, follow from the one-marking argument in
Section~\ref{sec:open}.

The proof of Theorem~\ref{thm:main} follows the induction in
\cite[Section~6.2]{CLP}.  The localization sequence separates degree-sixteen
homology into a normalized-boundary image and a lowest-weight Borel--Moore
quotient on $\M_{g,n}$.  Bergstr\"om--Faber--Payne vanishing makes the open
quotient zero outside a finite range.  Within that range, the published
Chow--K\"unneth generation results control every case except
$(3,12)$, $(5,8)$, $(6,6)$, and $(7,4)$.  Canning--Larson--Payne's
primitive-quotient formula removes one marking in these cases.  Virtual
cohomological dimension kills both the primitive local-system quotient and
the unmultiplied pullback term.  The remaining cotangent-line term comes
from a tautological group at the inclusive endpoint of their published
range, and Ionel's theorem makes that source zero.

On the boundary, every component of the normalization is a finite quotient
of either one smaller moduli stack or a product of two.  Odd--odd K\"unneth
summands in total degree sixteen vanish because one odd degree is at most
seven.  Positive even--even summands are Tate by the uniform result through
degree fourteen, and the two axis summands follow by induction.  Rational
averaging handles graph automorphisms, while semisimplicity of polarizable
pure Hodge structures splits the resulting same-weight extensions.

The same method isolates the remaining obstruction to the unrestricted
degree-sixteen problem.

\begin{theorem}\label{thm:reduction-intro}
Suppose that
\[
 W_{-16}H^{\BM}_{16}(\M_{8,n};\Q)
\]
is a direct sum of copies of $\Q(8)$ for $n=0,1$.  Then
$H^{16}(\Mbar_{g,n};\Q)$ is a direct sum of copies of $\Q(-8)$ for every
stable pair $(g,n)$.
\end{theorem}

Thus the genus-at-most-seven theorem is not an isolated collection of
examples: it leaves precisely two genus-eight open inputs.  The apparent
third endpoint, $\M_{8,2}$, vanishes by Wong's theorem
\cite[Theorem~4.1]{Wong}.  We return to this reduction in
Section~\ref{sec:consequences}.

\section{Mixed Hodge structures and the boundary sequence}
\label{sec:boundary}

\subsection{Weights, twists, and duality}

The Tate Hodge structure $\Q(m)$ has weight $-2m$ and Hodge type
$(-m,-m)$.  Thus $\Q(-8)$ has weight $16$ and type $(8,8)$, whereas
$\Q(8)$ has weight $-16$ and type $(-8,-8)$.  We use the increasing weight
filtration $W_\bullet$.  The mixed Hodge theory used below is due to Deligne
\cite{DeligneHodgeII,DeligneHodgeIII}; for cohomology and duality on
Deligne--Mumford stacks, see Behrend \cite{BehrendStacks}.  Throughout, the
rational cohomology of a Deligne--Mumford stack of finite type over $\C$
agrees, with its mixed Hodge structure, with that of its coarse moduli
space; for the stacks considered here the coarse space is a quasi-projective
variety with finite quotient singularities, so rational Poincar\'e duality
holds in the smooth case.

If $X$ is smooth of pure complex dimension $d$, Poincar\'e duality gives an
isomorphism of mixed Hodge structures
\begin{equation}\label{eq:pd-smooth}
 H_i^{\BM}(X;\Q)\cong H^{2d-i}(X;\Q)(d).
\end{equation}
If $X$ is proper, ordinary homology is the Hodge dual of ordinary
cohomology:
\begin{equation}\label{eq:proper-dual}
 H_i(X;\Q)\cong H^i(X;\Q)^\vee.
\end{equation}
For smooth proper $X$, these structures are polarizable and pure of weights
$i$ in cohomology and $-i$ in homology.

Every subobject or quotient of $\Q(m)^{\oplus N}$ in the category of
rational mixed Hodge structures is a direct sum of copies of $\Q(m)$.
Moreover, polarizable pure rational Hodge structures of fixed weight form a
semisimple category.  Consequently, an extension of copies of one Tate
structure by copies of the same Tate structure splits noncanonically.

\subsection{Normalization of the boundary}

Knudsen's clutching morphisms \cite{Knudsen} describe the normalized
boundary by one-edge stable graphs.  Here a stable graph of type $(g,n)$ is
a connected graph with legs labelled by $\{1,\ldots,n\}$ and with a
nonnegative integer $g_v$ at each vertex, such that
\[
 g=h^1(\Gamma)+\sum_v g_v,
 \qquad 2g_v-2+n_v>0
\]
for every vertex; $n_v$ counts incident half-edges and legs.  Its
automorphisms preserve every leg label and every vertex genus.  In
particular, the automorphism of a loop may exchange its two branches.
If $\Gamma$ has one edge, put
\[
 \Mbar_\Gamma:=\prod_{v\in V(\Gamma)}\Mbar_{g_v,n_v}.
\]
Then
\begin{equation}\label{eq:boundary-graphs}
 \widetilde{\partial\M}_{g,n}
 \cong
 \coprod_{\substack{[\Gamma]\text{ stable of type }(g,n)\\
                     |E(\Gamma)|=1}}
 \left[\Mbar_\Gamma/\Aut(\Gamma)\right].
\end{equation}
This is the description used in \cite[Section~2.1]{CLP}.  A loop gives the
nonseparating component $[\Mbar_{g-1,n+2}/S_2]$; an edge joining two
vertices gives a finite quotient of a product.

For a finite group $G$ acting algebraically on a complex Deligne--Mumford
stack $Y$, the Cartan--Leray spectral sequence and rational averaging give
\begin{equation}\label{eq:quotient-invariants}
 H^*([Y/G];\Q)\cong H^*(Y;\Q)^G
\end{equation}
as mixed Hodge structures.  Indeed, $H^p(G,V)=0$ for $p>0$ over $\Q$.
Thus finite graph-automorphism quotients do not introduce new Hodge types.

\subsection{The right-exact sequence}

Let
\[
 \nu\colon\widetilde{\partial\M}_{g,n}\longrightarrow\Mbar_{g,n}
\]
be the natural finite map.

\begin{proposition}[Canning--Larson--Payne]\label{prop:boundary-sequence}
For every stable pair $(g,n)$ and every integer $i$, there is a right-exact
sequence of mixed Hodge structures
\begin{equation}\label{eq:boundary-sequence}
 H_i(\widetilde{\partial\M}_{g,n};\Q)
 \xrightarrow{\nu_*}
 H_i(\Mbar_{g,n};\Q)
 \longrightarrow
 W_{-i}H_i^{\BM}(\M_{g,n};\Q)
 \longrightarrow0.
\end{equation}
\end{proposition}

This is precisely Equation~(6.1) of \cite{CLP}.  We use that published
mixed-Hodge-theoretic statement directly; in particular, the first map in
\eqref{eq:boundary-sequence} has neither a degree shift nor a Tate twist.

For $i=16$, let $I_{g,n}$ be the image of the first map.  We obtain a short
exact sequence
\begin{equation}\label{eq:boundary-short}
 0\longrightarrow I_{g,n}
 \longrightarrow H_{16}(\Mbar_{g,n};\Q)
 \longrightarrow W_{-16}H_{16}^{\BM}(\M_{g,n};\Q)
 \longrightarrow0.
\end{equation}

\section{The open contribution in genus at most seven}
\label{sec:open}

\subsection{The primitive quotient}

For $1\leq i\leq n$, let
\[
 \pi_i\colon\M_{g,n}\longrightarrow\M_{g,n-1}
\]
forget the $i$th marking, and let $\psi_i\in H^2(\M_{g,n};\Q)$ be the first
Chern class of its cotangent line.  Following \cite[Sections~2 and 3]{CLP},
define
\begin{align}
 \Phi^k_{g,n}
   &:=\sum_{i=1}^n\pi_i^*W_kH^k(\M_{g,n-1};\Q),
   \label{eq:Phi}\\
 \Psi^k_{g,n}
   &:=\sum_{i=1}^n\psi_i\smile\Phi^{k-2}_{g,n}.
   \label{eq:Psi}
\end{align}
Both are subspaces of $W_kH^k(\M_{g,n};\Q)$.  Let
$\lambda=(\lambda_1\geq\cdots\geq\lambda_g\geq0)$ be a partition with
$|\lambda|:=\sum_i\lambda_i=n$.  Write $\mathbb V_\lambda$ for the
associated irreducible symplectic local system on $\M_g$, and write
$V_{\lambda^{\mathsf T}}$ for the irreducible $S_n$-representation
corresponding to the transpose partition.  These are the conventions of
\cite[Section~3]{CLP}.

\begin{proposition}[Canning--Larson--Payne]\label{prop:primitive}
If $g\geq2$, there is an isomorphism of $\Q$-vector spaces
\begin{equation}\label{eq:primitive}
 \frac{W_kH^k(\M_{g,n};\Q)}
      {\Phi^k_{g,n}+\Psi^k_{g,n}}
 \cong
 \bigoplus_{|\lambda|=n}
 W_kH^{k-n}(\M_g;\mathbb V_\lambda)
 \otimes V_{\lambda^{\mathsf T}}.
\end{equation}
\end{proposition}

\begin{proof}
This is \cite[Lemma~3.1(a)]{CLP}.  In its proof, the summands of relative
cohomological degree two are identified by the relative projector $\pi_2$
of Petersen--Tavakol--Yin \cite[Sections~5.1 and 5.2.2]{PTY}.
\end{proof}

We use only the following consequence: if every summand on the right of
\eqref{eq:primitive} vanishes, then
$W_kH^k(\M_{g,n};\Q)=\Phi^k_{g,n}+\Psi^k_{g,n}$, so the vanishing of
$\Phi^k_{g,n}$ and of $\Psi^k_{g,n}$ forces $W_kH^k(\M_{g,n};\Q)=0$.  No
compatibility of \eqref{eq:primitive} with Hodge structures is needed.

The following coefficient form of virtual-cohomological-dimension
vanishing will be used repeatedly.

\begin{lemma}\label{lem:vcd-rational}
Let $\Gamma$ be virtually torsion-free of virtual cohomological dimension
$d$, and let $L$ be a finite-dimensional rational $\Gamma$-module.  Then
\[
 H^q(\Gamma;L)=0\qquad(q>d).
\]
\end{lemma}

\begin{proof}
Choose a torsion-free normal subgroup $\Gamma'\triangleleft\Gamma$ of finite
index; by Serre's theorem it has cohomological dimension $d$.  Then
$H^q(\Gamma';L)=0$ for $q>d$.
In the Hochschild--Serre spectral sequence for $\Gamma'$ in $\Gamma$, the
quotient is finite.  Its higher rational cohomology vanishes, so
\[
 H^q(\Gamma;L)\cong H^q(\Gamma';L)^{\Gamma/\Gamma'}=0
 \qquad(q>d).
\]
See \cite[Chapters~VII and~VIII]{Brown}.
\end{proof}

Teichm\"uller space is contractible, so the relevant stack cohomology with
local coefficients agrees with mapping-class-group cohomology.  Harer proved
that, for $g\geq2$ and $n>0$,
\begin{equation}\label{eq:harer-vcd}
 \vcd(\Mod_g)=4g-5,
 \qquad
 \vcd(\PMod_{g,n})=4g-4+n
\end{equation}
\cite[Theorem~4.1]{Harer}.

\subsection{Chow--K\"unneth generation and vanishing}

A Chow group in this paper always has rational coefficients:
$A^*(X)=A^*(X)_{\Q}$.  Let $R^*(\Mbar_{g,n})\subset A^*(\Mbar_{g,n})$ be the
tautological ring, and let $R^*(\M_{g,n})\subset A^*(\M_{g,n})$ be its image
under restriction; since all boundary generators restrict to zero, the ring
$R^*(\M_{g,n})$ is generated by the $\psi$- and $\kappa$-classes.  Define
tautological cohomology by
\begin{equation}\label{eq:RH-definition}
 \begin{aligned}
 \RH^{2c}(\M_{g,n};\Q)
 &:=\im\!\left(R^c(\M_{g,n})\xrightarrow{\mathrm{cl}}
 H^{2c}(\M_{g,n};\Q)\right),\\
 \RH^{2c+1}(\M_{g,n};\Q)&:=0.
 \end{aligned}
\end{equation}

A smooth algebraic stack $X$ has the \emph{Chow--K\"unneth generation
property} if the exterior product
\[
 A^*(X)\otimes A^*(Y)\longrightarrow A^*(X\times Y)
\]
is surjective for every finite-type stack $Y$ admitting a stratification by
global quotient stacks.  Canning--Larson--Payne prove that if a smooth open
substack of a smooth proper Deligne--Mumford stack has this property, then
its cycle class map surjects onto all lowest-weight cohomology
\cite[Lemma~4.3]{CLP}.

Their Table~1 gives marking bounds
\begin{equation}\label{eq:c-table}
\begin{array}{c|cccccccc}
g&0&1&2&3&4&5&6&7\\ \hline
c(g)&\infty&10&10&11&11&7&5&3
\end{array}
\end{equation}
such that, for $g\leq7$ and $n\leq c(g)$,
\begin{equation}\label{eq:CLP-pure-taut}
 W_kH^k(\M_{g,n};\Q)=\RH^k(\M_{g,n};\Q)
\end{equation}
for every $k$ \cite[Proposition~4.5]{CLP}.  The geometric results behind most
of these bounds come from Canning--Larson \cite{CanningLarsonStable}; the
direct locators used here are Table~1, Lemma~4.3, and Proposition~4.5 of
\cite{CLP}.

The tautological vanishing needed below combines the pointed theorem of
Ionel with the unpointed theorem of Looijenga.

\begin{proposition}[Ionel--Looijenga]\label{prop:ionel}
For $g\geq2$, tautological cohomology on the open moduli stack satisfies
\begin{equation}\label{eq:ionel}
 \RH^{2c}(\M_{g,n};\Q)=0
 \quad\text{if}\quad
 \begin{cases}
 c\geq g,&n>0,\\
 c\geq g-1,&n=0.
 \end{cases}
\end{equation}
\end{proposition}

\begin{proof}
For $n>0$, Ionel proves that every product of total degree at least $g$ in
descendant and tautological classes vanishes in the rational cohomology of
$\M_{g,n}$ \cite[Theorem~0.1]{Ionel}.  Since $R^*(\M_{g,n})$ is generated
by the $\psi$- and $\kappa$-classes, every class in $R^c(\M_{g,n})$ with
$c\geq g$ is a combination of such products, so its cohomology class
vanishes.  For $n=0$, Looijenga proves that the rational
tautological Chow ring of $\M_g$ vanishes in codimension greater than
$g-2$ \cite[Theorem~1.2]{Looijenga}.  Passing through the cycle map and
using \eqref{eq:RH-definition} gives \eqref{eq:ionel}.
\end{proof}

\subsection{The critical line}

Put $d=3g-3+n$ and $k=2d-16$.  For $2g+n=18$, the number of markings
$n=18-2g$ is positive in every genus considered here, and
\begin{equation}\label{eq:critical-arithmetic}
 k=4g-4+n=\vcd(\PMod_{g,n}).
\end{equation}
The relevant cases in genera three through seven are displayed below.

\begin{table}[ht]
\centering
\caption{The critical open cases in genus at most seven.}
\label{tab:critical}
\small
\begin{tabular}{@{}cccc@{}}
\toprule
$(g,n)$&$d$&$k=2d-16$&source of the open calculation\\
\midrule
$(3,12)$&18&20&$n-1=c(3)$; one-marking argument\\
$(4,10)$&19&22&\cite[Proposition~4.5]{CLP}\\
$(5,8)$ &20&24&$n-1=c(5)$; one-marking argument\\
$(6,6)$ &21&26&$n-1=c(6)$; one-marking argument\\
$(7,4)$ &22&28&$n-1=c(7)$; one-marking argument\\
\bottomrule
\end{tabular}
\end{table}

\begin{proposition}\label{prop:endpoint-vanishing}
For
\[
 (g,n,k)=(3,12,20),\ (5,8,24),\ (6,6,26),\ (7,4,28),
\]
one has
\begin{equation}\label{eq:endpoint-zero}
 W_kH^k(\M_{g,n};\Q)=0.
\end{equation}
\end{proposition}

\begin{proof}
Fix one of the four triples and apply
Proposition~\ref{prop:primitive}.  Since $k=4g-4+n$, the degree on the
right of \eqref{eq:primitive} is
\[
 k-n=4g-4>4g-5=\vcd(\Mod_g).
\]
Lemma~\ref{lem:vcd-rational} makes every local-system summand vanish.
Hence
\begin{equation}\label{eq:PhiPsi}
 W_kH^k(\M_{g,n};\Q)=\Phi^k_{g,n}+\Psi^k_{g,n}.
\end{equation}
Moreover,
\[
 \vcd(\PMod_{g,n-1})=4g-4+(n-1)=k-1,
\]
so $H^k(\M_{g,n-1};\Q)=0$.  Thus $\Phi^k_{g,n}=0$.

It remains to consider the source of $\Psi^k_{g,n}$.  In all four cases,
$n-1=c(g)$.  By \cite[Table~1, Lemma~4.3, and
Proposition~4.5]{CLP},
\[
 W_{k-2}H^{k-2}(\M_{g,n-1};\Q)
 =\RH^{k-2}(\M_{g,n-1};\Q).
\]
The four Chow codimensions $(k-2)/2$ are respectively $9$, $11$, $12$, and
$13$, each at least $g$.  Proposition~\ref{prop:ionel} therefore makes
these source groups zero.  Equations \eqref{eq:Psi} and \eqref{eq:PhiPsi}
now imply \eqref{eq:endpoint-zero}.
\end{proof}

\begin{corollary}\label{cor:critical-vanishing}
All five groups in \eqref{eq:critical-intro} vanish.  Equivalently, for
$3\leq g\leq7$ and $2g+n=18$,
\begin{equation}\label{eq:critical-bm-zero}
 W_{-16}H^{\BM}_{16}(\M_{g,n};\Q)=0.
\end{equation}
\end{corollary}

\begin{proof}
Proposition~\ref{prop:endpoint-vanishing} handles $(3,12)$, $(5,8)$,
$(6,6)$, and $(7,4)$.  At $(4,10)$, Proposition~4.5 of \cite{CLP} identifies the pure
weight group with tautological cohomology; its codimension is $11\geq4$, so
Proposition~\ref{prop:ionel} makes it zero.

For every row, $k=2d-16$.  Poincar\'e duality \eqref{eq:pd-smooth} gives
\[
 W_{-16}H^{\BM}_{16}(\M_{g,n};\Q)
 \cong W_kH^k(\M_{g,n};\Q)(d),
\]
which proves \eqref{eq:critical-bm-zero}.
\end{proof}

\begin{remark}\label{rem:liu}
A recent preprint of Liu \cite{Liu} announces the stronger Chow-theoretic
statements that $A^*(\M_{5,8})$ and $A^*(\M_{5,9})$ are tautological and
that these stacks have the Chow--K\"unneth generation property, by a direct
analysis of configurations of points on canonical genus-five curves.  Those
results would give the $(5,8)$ row of \eqref{eq:critical-intro} by the same
mechanism as the $(4,10)$ row.  The argument given here is independent of
that preprint: the $(5,8)$ case of Proposition~\ref{prop:endpoint-vanishing}
uses only the published range $c(5)=7$ of \cite{CLP,CanningLarsonStable},
together with Proposition~\ref{prop:primitive}, Harer's theorem, and Ionel's
theorem.
\end{remark}

\subsection{The open term throughout the range}

\begin{proposition}\label{prop:open-all}
Let $3\leq g\leq7$ and let $(g,n)$ be stable.  Then, for some
$a_{g,n}\geq0$,
\begin{equation}\label{eq:open-all}
 W_{-16}H_{16}^{\BM}(\M_{g,n};\Q)
 \cong\Q(8)^{\oplus a_{g,n}}.
\end{equation}
Moreover, $a_{g,n}=0$ whenever $2g+n\geq11$.
\end{proposition}

\begin{proof}
If $n\geq2$ and $2g+n>18$, the Borel--Moore group itself vanishes by
\cite[Proposition~2.1]{BFP}; indeed, that proposition gives
\[
 H_i^{\BM}(\M_{g,n};\Q)=0\qquad(i<2g-2+n).
\]
We may therefore assume either $n\leq1$ or $2g+n\leq18$.

Put $d=3g-3+n$ and $j=2d-16$.  If $j<0$, the group is zero by
\eqref{eq:pd-smooth}.  Assume $j\geq0$.  The remaining marking bounds are
\[
\begin{array}{c|ccccc}
g&3&4&5&6&7\\ \hline
\text{largest }n&12&10&8&6&4.
\end{array}
\]
For $g=3,5,6,7$, every pair with $n$ strictly below the listed bound lies
in the ranges \eqref{eq:c-table}, and for $g=4$ every pair through
$(4,10)$ does.  The only pairs outside these direct ranges are
$(3,12)$, $(5,8)$, $(6,6)$, and $(7,4)$.

In the direct Canning--Larson--Payne range,
\eqref{eq:CLP-pure-taut} shows that $W_jH^j$ is algebraic and hence a direct
sum of $\Q(-j/2)$: here $j$ is even, and every codimension-$j/2$ cycle class
has Hodge type $(j/2,j/2)$.  The four exceptional groups vanish by
Proposition~\ref{prop:endpoint-vanishing}.  Therefore, in every remaining
case,
\[
 W_jH^j(\M_{g,n};\Q)
 \cong\Q(-j/2)^{\oplus a_{g,n}}
\]
for some $a_{g,n}\geq0$.  Since $j/2=d-8$, Poincar\'e duality gives
\begin{align*}
 W_{-16}H_{16}^{\BM}(\M_{g,n};\Q)
 &\cong W_jH^j(\M_{g,n};\Q)(d)\\
 &\cong\Q(-j/2)(d)^{\oplus a_{g,n}}
  =\Q(8)^{\oplus a_{g,n}}.
\end{align*}

Suppose finally that $2g+n\geq11$.  If $n>0$, then
\[
 \frac j2=3g+n-11\geq g,
\]
so Proposition~\ref{prop:ionel} makes every direct-range group above zero.
If $n=0$, the inequality forces $g\geq6$, and
$j/2=3g-11\geq g-1$, so the unpointed clause of the same proposition
applies.  The four exceptional groups and the cases removed at the start
of the proof already vanish.  Hence $a_{g,n}=0$ in the asserted range.
\end{proof}

\section{Compact low-degree inputs and products}
\label{sec:compact-inputs}

We collect the compact results used in the boundary induction.

\begin{proposition}\label{prop:compact-inputs}
The following statements hold for stable pairs.
\begin{enumerate}[label=\textup{(\roman*)}]
\item $H^q(\Mbar_{g,n};\Q)=0$ for odd $q\leq9$.
\item If $0\leq q\leq14$ is even, then
\[
 H^q(\Mbar_{g,n};\Q)\cong\Q(-q/2)^{\oplus b_q}
\]
for some $b_q\geq0$.
\item For every $n$, the group $H^{16}(\Mbar_{0,n};\Q)$ is a direct sum of
copies of $\Q(-8)$ whenever $(0,n)$ is stable.
\item For every stable $(1,n)$ and $(2,n)$,
\[
 H^{16}(\Mbar_{g,n};\Q)\cong\Q(-8)^{\oplus b_{g,n}}.
\]
\end{enumerate}
\end{proposition}

\begin{proof}
Part~(i) is the theorem of Bergstr\"om--Faber--Payne
\cite[Theorem~1.1]{BFP}.  For part~(ii), recall that
Canning--Larson--Payne write
\[
 H_q(\Mbar_{g,n}):=H^{2d-q}(\Mbar_{g,n}),
 \qquad d=3g-3+n.
\]
Their Theorem~1.5(3) says that this complementary cohomology is
tautological for even $q\leq14$.  Hence it is a sum of copies of
$\Q(-(d-q/2))$.  Poincar\'e duality in the form
\[
 H^q(\Mbar_{g,n})
 \cong H^{2d-q}(\Mbar_{g,n})^\vee(-d)
\]
gives the Tate type $\Q(-q/2)$.  This transfers Hodge type, not
tautological generation in degree $q$.

Keel proves that the cycle class map from the Chow ring of
$\Mbar_{0,n}$ to cohomology is an isomorphism \cite{Keel}, which gives
part~(iii).  Petersen proves that the even cohomology of
$\Mbar_{1,n}$ is spanned by algebraic strata classes and is therefore
tautological \cite[Corollary~1.2]{PetersenGenusOne}.
Finally, \cite[Corollary~2.7]{CLP} states directly that
\[
 H^{2k}(\Mbar_{2,n};\Q)=\RH^{2k}(\Mbar_{2,n};\Q)
 \qquad(k\leq10).
\]
Taking $k=8$ proves part~(iv).
\end{proof}

\begin{lemma}\label{lem:kunneth}
Let $(g_1,n_1)$ and $(g_2,n_2)$ be stable.  If
$H^{16}(\Mbar_{g_i,n_i};\Q)$ is a direct sum of copies of $\Q(-8)$ for
$i=1,2$, then
\begin{equation}\label{eq:kunneth}
 H^{16}(\Mbar_{g_1,n_1}\times\Mbar_{g_2,n_2};\Q)
 \cong\Q(-8)^{\oplus N}
\end{equation}
for some $N\geq0$.
\end{lemma}

\begin{proof}
The K\"unneth isomorphism is compatible with Hodge structures:
\[
 H^{16}(X\times Y;\Q)
 \cong\bigoplus_{p+q=16}H^p(X;\Q)\otimes H^q(Y;\Q).
\]
If $p$ and $q$ are odd, then $\min(p,q)\leq7$, so one factor vanishes by
Proposition~\ref{prop:compact-inputs}(i).  If $p,q>0$ are even, then both
are at most fourteen, and part~(ii) identifies their tensor product with a
sum of
\[
 \Q(-p/2)\otimes\Q(-q/2)=\Q(-8).
\]
The two remaining summands have bidegrees $(16,0)$ and $(0,16)$ and are
covered by the hypotheses, since every moduli stack in question is
connected and has $H^0\cong\Q(0)$.
\end{proof}

\section{Proof in genus at most seven}
\label{sec:induction}

\begin{proof}[Proof of Theorem~\ref{thm:main}]
We argue by induction on $d=3g-3+n$, simultaneously for all stable pairs
with $g\leq7$.  The cases $g=0,1,2$ hold for every $n$ by
Proposition~\ref{prop:compact-inputs}(iii)--(iv), so they may be used at any
stage of the induction.  Assume $3\leq g\leq7$.

Proposition~\ref{prop:open-all} gives
\begin{equation}\label{eq:induction-open}
 W_{-16}H_{16}^{\BM}(\M_{g,n};\Q)
 \cong\Q(8)^{\oplus a_{g,n}}.
\end{equation}
Consider a one-edge stable graph $\Gamma$ of type $(g,n)$.  Every vertex
has genus at most $g$, and partial normalization at the distinguished node
lowers the total complex dimension by one.  More explicitly, a loop has
ordered atlas $\Mbar_{g-1,n+2}$ of dimension $d-1$, while a separating edge
has ordered atlas
\[
 \Mbar_{g_1,n_1}\times\Mbar_{g_2,n_2},
 \qquad
 \dim\Mbar_{g_1,n_1}+\dim\Mbar_{g_2,n_2}=d-1.
\]
Each individual factor has dimension smaller than $d$.

For a loop, the induction hypothesis applies to the single factor.  For a
separating edge, it applies to both factors, and
Lemma~\ref{lem:kunneth} gives the required type for their product.  Taking
the finite graph-automorphism invariants in
\eqref{eq:quotient-invariants} preserves this type.  Summing over the
finitely many one-edge graphs in \eqref{eq:boundary-graphs}, we obtain
\begin{align}
 H^{16}(\widetilde{\partial\M}_{g,n};\Q)
   &\cong\Q(-8)^{\oplus N_{g,n}},
   \label{eq:boundary-cohom-induction}\\
 H_{16}(\widetilde{\partial\M}_{g,n};\Q)
   &\cong\Q(8)^{\oplus N_{g,n}}.
   \label{eq:boundary-hom-induction}
\end{align}

The boundary image $I_{g,n}$ in \eqref{eq:boundary-short} is a quotient of
\eqref{eq:boundary-hom-induction}.  Hence
$I_{g,n}\cong\Q(8)^{\oplus e_{g,n}}$ for some $e_{g,n}$.  Combining this
with \eqref{eq:induction-open}, the boundary sequence becomes
\begin{equation}\label{eq:induction-extension}
 0\longrightarrow\Q(8)^{\oplus e_{g,n}}
 \longrightarrow H_{16}(\Mbar_{g,n};\Q)
 \longrightarrow\Q(8)^{\oplus a_{g,n}}
 \longrightarrow0.
\end{equation}
The middle term is the dual of a polarizable pure Hodge structure of weight
$16$, so it is polarizable and pure of weight $-16$.  Semisimplicity splits
\eqref{eq:induction-extension} and gives
\[
 H_{16}(\Mbar_{g,n};\Q)
 \cong\Q(8)^{\oplus(a_{g,n}+e_{g,n})}.
\]
Dualizing proves \eqref{eq:main}.  Every use of the induction hypothesis was
at strictly smaller dimension, so the induction is well founded.
\end{proof}

\section{The six-pointed genus-six case and the genus-eight reduction}
\label{sec:consequences}

\subsection{Boundary surjectivity in the genus-six case}

The dimension calculation is
\[
 \dim_{\C}\Mbar_{6,6}=3\cdot6-3+6=21.
\]
The nonseparating component of its normalized boundary is
\begin{equation}\label{eq:nonsep-66}
 [\Mbar_{5,8}/S_2],
\end{equation}
where $S_2$ exchanges the two branches of the distinguished node.  A
separating component is a finite quotient of
\begin{equation}\label{eq:sep-66}
 \Mbar_{g_1,n_1}\times\Mbar_{g_2,n_2},
 \qquad g_1+g_2=6,
 \quad n_1+n_2=8.
\end{equation}
All vertex genera are at most six.  Theorem~\ref{thm:main},
Lemma~\ref{lem:kunneth}, and rational invariants therefore show that, for
some $N\geq0$,
\begin{equation}\label{eq:boundary-66-type}
 H^{16}(\widetilde{\partial\M}_{6,6};\Q)
 \cong\Q(-8)^{\oplus N},
 \qquad
 H_{16}(\widetilde{\partial\M}_{6,6};\Q)
 \cong\Q(8)^{\oplus N}.
\end{equation}

\begin{proof}[Proof of Corollary~\ref{cor:66-intro}]
Corollary~\ref{cor:critical-vanishing} gives
\[
 W_{-16}H^{\BM}_{16}(\M_{6,6};\Q)=0.
\]
The right-exact sequence \eqref{eq:boundary-sequence} is therefore exactly
the surjection \eqref{eq:66-boundary-surj-intro}.  Its target is a quotient
of the second group in \eqref{eq:boundary-66-type}, so
\[
 H_{16}(\Mbar_{6,6};\Q)\cong\Q(8)^{\oplus b_{16}}.
\]
Proper duality gives the asserted cohomological statement.
\end{proof}

The conclusion concerns the image of boundary \emph{homology}.  It does not
assert that $H^{16}(\Mbar_{6,6})$ is generated by algebraic boundary cycles.

\subsection{Reduction to two genus-eight inputs}

The Borel--Moore vanishing used in Proposition~\ref{prop:open-all} holds in
all genera:
\begin{equation}\label{eq:BFP-BM}
 H_{16}^{\BM}(\M_{g,n};\Q)=0
 \quad\text{if}\quad
 \begin{cases}
 16<2g,&n=0,1,\\
 16<2g-2+n,&n\geq2.
 \end{cases}
\end{equation}
This is \cite[Proposition~2.1]{BFP} together with duality between
Borel--Moore homology and compactly supported cohomology.

The equality case left open by \eqref{eq:BFP-BM} at $(g,n)=(8,2)$ is
settled by Wong:
\begin{equation}\label{eq:wong-82}
 H^{30}(\M_{8,2};\Q)=0,
 \qquad H_{16}^{\BM}(\M_{8,2};\Q)=0.
\end{equation}
Indeed, the first equality is the specialization to $g=8$ of
\cite[Theorem~4.1]{Wong}, which states
$H^{4g-2}(\M_{g,2};\Q)=0$ for every $g\geq1$.  Since
$\dim_{\C}\M_{8,2}=23$, the second equality follows from
\eqref{eq:pd-smooth}.

\begin{proof}[Proof of Theorem~\ref{thm:reduction-intro}]
By Theorem~\ref{thm:main}, all vertex types of genus at most seven are
already known.  We first use the two assumed open groups to establish the
first compact genus-eight base cases.  For $(8,0)$, every vertex of a one-edge
stable graph has genus at most seven: a genus-eight vertex on a separating
edge would leave a vertex of type $(0,1)$, which is unstable, and the loop
atlas is $\Mbar_{7,2}$.  The boundary argument from
Section~\ref{sec:induction}, together with the assumed open quotient, proves
the compact assertion.  The same is true for $(8,1)$: a rational component
carrying the unique original marking and one branch of the node would be
unstable, so no genus-eight vertex occurs in its normalized boundary.

For $(8,2)$, the open term vanishes by \eqref{eq:wong-82}.  The only new
possibility in its normalized boundary with a genus-eight vertex is a
rational tail carrying the two original markings; the other factor is then
of type $(8,1)$, which has just been established.  The boundary sequence
therefore proves the compact assertion for $(8,2)$ without a third
genus-eight hypothesis.

We now prove the assertion for all remaining stable pairs by induction on
$d=3g-3+n$.

For $g=8$ and $n\geq3$, the second inequality in \eqref{eq:BFP-BM} holds,
while $n=2$ is covered by \eqref{eq:wong-82}.
For $g\geq9$, it holds for every $n\geq2$, and the first inequality holds
for $n=0,1$.  Thus the open quotient in \eqref{eq:boundary-short} is zero
outside the two assumed genus-eight pairs.

Every vertex of a one-edge graph has smaller dimension than the original
pair.  The induction hypothesis, Lemma~\ref{lem:kunneth}, and finite-group
invariants therefore give Tate type $\Q(-8)$ for the degree-sixteen
cohomology of the normalized boundary.  Its homology is of type $\Q(8)$,
so the boundary sequence and semisimplicity prove the assertion for the
original pair.  This completes the induction.
\end{proof}

For clarity, the two remaining open groups correspond under
\eqref{eq:pd-smooth} to the following ordinary lowest-weight groups:
\begin{equation}\label{eq:genus8-table}
\begin{array}{c|ccc}
n&d=21+n&j=2d-16&j/2\\ \hline
0&21&26&13\\
1&22&28&14\\
\end{array}
\end{equation}
If $\M_{8,n}$ had the Chow--K\"unneth generation property and tautological
rational Chow ring for $n=0,1$, Lemma~4.3 of \cite{CLP} would identify
$W_jH^j$ with tautological cohomology.  Proposition~\ref{prop:ionel} would
then make both groups zero: the codimensions are $13$ and $14$, whereas the
respective vanishing thresholds are $7$ and $8$.  Canning and
Larson have proved that the rational Chow ring of $\M_8$ is tautological
\cite{CanningLarson789}; the additional Chow--K\"unneth input is a separate
requirement, since tautologicality of the Chow ring alone does not control
lowest-weight cohomology.

\begin{remark}\label{rem:scope}
All results are over $\C$, use rational coefficients, and concern
Deligne--Mumford stacks rather than their coarse moduli spaces.  No integral,
$\ell$-adic, point-counting, or exact-rank statement is asserted.
\end{remark}

\section*{Acknowledgments}

The principal proof development, mathematical analysis, drafting, and review
underlying this manuscript were carried out by the artificial intelligence
collaborators named in the byline.

\begin{thebibliography}{99}

\bibitem{BehrendStacks}
K.~Behrend,
\emph{Cohomology of stacks},
in \emph{Intersection theory and moduli}, ICTP Lecture Notes, vol.~XIX,
Abdus Salam International Centre for Theoretical Physics, Trieste, 2004,
pp.~249--294.

\bibitem{BFP}
J.~Bergstr\"om, C.~Faber, and S.~Payne,
\emph{Polynomial point counts and odd cohomology vanishing on moduli spaces
of stable curves},
Ann. of Math. (2) \textbf{199} (2024), no.~3, 1323--1365,
\href{https://doi.org/10.4007/annals.2024.199.3.7}
{doi:10.4007/annals.2024.199.3.7}.

\bibitem{Brown}
K.~S. Brown,
\emph{Cohomology of groups},
Graduate Texts in Mathematics, vol.~87, Springer-Verlag, New York, 1982,
\href{https://doi.org/10.1007/978-1-4684-9327-6}
{doi:10.1007/978-1-4684-9327-6}.

\bibitem{CanningLarson789}
S.~Canning and H.~Larson,
\emph{The Chow rings of the moduli spaces of curves of genus $7$, $8$, and
$9$},
J. Algebraic Geom. \textbf{33} (2024), no.~1, 55--116,
\href{https://doi.org/10.1090/jag/818}{doi:10.1090/jag/818}.

\bibitem{CanningLarsonStable}
S.~Canning and H.~Larson,
\emph{On the Chow and cohomology rings of moduli spaces of stable curves},
J. Eur. Math. Soc. \textbf{28} (2026), no.~9, 3869--3918,
\href{https://doi.org/10.4171/JEMS/1543}{doi:10.4171/JEMS/1543}.

\bibitem{CLP}
S.~Canning, H.~Larson, and S.~Payne,
\emph{Extensions of tautological rings and motivic structures in the
cohomology of $\overline{\mathcal M}_{g,n}$},
Forum Math. Pi \textbf{12} (2024), article no.~e23,
\href{https://doi.org/10.1017/fmp.2024.24}{doi:10.1017/fmp.2024.24}.

\bibitem{DeligneHodgeII}
P.~Deligne,
\emph{Th\'eorie de Hodge. II},
Publ. Math. Inst. Hautes \'Etudes Sci. \textbf{40} (1971), 5--57,
\href{https://doi.org/10.1007/BF02684692}{doi:10.1007/BF02684692}.

\bibitem{DeligneHodgeIII}
P.~Deligne,
\emph{Th\'eorie de Hodge. III},
Publ. Math. Inst. Hautes \'Etudes Sci. \textbf{44} (1974), 5--77,
\href{https://doi.org/10.1007/BF02685881}{doi:10.1007/BF02685881}.

\bibitem{Harer}
J.~L. Harer,
\emph{The virtual cohomological dimension of the mapping class group of an
orientable surface},
Invent. Math. \textbf{84} (1986), no.~1, 157--176,
\href{https://doi.org/10.1007/BF01388737}{doi:10.1007/BF01388737}.

\bibitem{Ionel}
E.-N. Ionel,
\emph{Topological recursive relations in $H^{2g}(\mathcal M_{g,n})$},
Invent. Math. \textbf{148} (2002), no.~3, 627--658,
\href{https://doi.org/10.1007/s002220100205}
{doi:10.1007/s002220100205}.

\bibitem{Keel}
S.~Keel,
\emph{Intersection theory of moduli space of stable $n$-pointed curves of
genus zero},
Trans. Amer. Math. Soc. \textbf{330} (1992), no.~2, 545--574,
\href{https://doi.org/10.1090/S0002-9947-1992-1034665-0}
{doi:10.1090/S0002-9947-1992-1034665-0}.

\bibitem{Knudsen}
F.~F. Knudsen,
\emph{The projectivity of the moduli space of stable curves. II. The stacks
$\overline{\mathcal M}_{g,n}$},
Math. Scand. \textbf{52} (1983), 161--199,
\href{https://doi.org/10.7146/math.scand.a-12001}
{doi:10.7146/math.scand.a-12001}.

\bibitem{Liu}
Y.~Liu,
\emph{On the Chow rings of the moduli spaces $\mathcal M_{5,8}$ and
$\mathcal M_{5,9}$},
preprint, 2025,
\href{https://arxiv.org/abs/2509.02950}{arXiv:2509.02950}.

\bibitem{Looijenga}
E.~Looijenga,
\emph{On the tautological ring of $\mathcal M_g$},
Invent. Math. \textbf{121} (1995), no.~2, 411--419,
\href{https://doi.org/10.1007/BF01884306}
{doi:10.1007/BF01884306}.

\bibitem{PayneSurvey}
S.~Payne,
\emph{On the Hodge and Tate conjectures for moduli spaces of curves},
preprint, 2026,
\href{https://arxiv.org/abs/2605.20453}{arXiv:2605.20453}.

\bibitem{PetersenGenusOne}
D.~Petersen,
\emph{The structure of the tautological ring in genus one},
Duke Math. J. \textbf{163} (2014), no.~4, 777--793,
\href{https://doi.org/10.1215/00127094-2429916}
{doi:10.1215/00127094-2429916}.

\bibitem{PTY}
D.~Petersen, M.~Tavakol, and Q.~Yin,
\emph{Tautological classes with twisted coefficients},
Ann. Sci. \`Ec. Norm. Sup\'er. (4) \textbf{54} (2021), no.~5, 1179--1236,
\href{https://doi.org/10.24033/asens.2479}{doi:10.24033/asens.2479}.

\bibitem{Wong}
Y.~M. Wong,
\emph{An algorithm to compute the fundamental classes of spin components
of strata of differentials},
Int. Math. Res. Not. IMRN (2024), no.~6, 4893--4962,
\href{https://doi.org/10.1093/imrn/rnad208}
{doi:10.1093/imrn/rnad208}.

\end{thebibliography}

\end{document}
